
\chapter*{Utiliser les corrigés}
\addcontentsline{toc}{chapter}{Utiliser les corrigés}
Les 180 exercices et les huit problèmes sont corrigés dans l'ordre du recueil.
Les solutions expliquent les choix de méthode et les conditions de validité.
Une autre rédaction correcte est possible ; l'essentiel est de justifier les
étapes et de répondre à la question posée.

Avant de lire, noter la tentative, l'endroit du blocage et une éventuelle
conjecture. Après lecture, fermer le corrigé et reconstruire le raisonnement.
Classer l'erreur : définition, domaine, calcul, hypothèse, logique, unité ou
interprétation. Cette analyse permet de choisir la révision utile.

Les résultats d'un exemple numérique ne démontrent pas une propriété générale.
Inversement, une démonstration ne dispense pas de vérifier que les données
d'un exercice remplissent ses hypothèses.
